\documentclass[12pt, twoside]{article}
%\documentclass[12pt, twoside]{article}
\usepackage[letterpaper, margin=1in, headsep=0.2in]{geometry}
\setlength{\headheight}{0.6in}
%\usepackage[english]{babel}
\usepackage[utf8]{inputenc}
\usepackage{microtype}
\usepackage{amsmath}
\usepackage{amssymb}
%\usepackage{amsfonts}
\usepackage[nomessages]{fp} %\FPeval{\var-name}{2*sin(pi/6)}
\usepackage{siunitx} %units in math. eg 20\milli\meter
\usepackage{yhmath} % for arcs, overparenth command
\usepackage{tikz} %graphics
\usetikzlibrary{quotes, angles, arrows, arrows.meta}
\usepackage{pgfplots}
\usepgfplotslibrary{fillbetween}
\pgfplotsset{compat=1.16}
\usepackage{graphicx} %consider setting \graphicspath{{images/}}
\usepackage{parskip} %no paragraph indent
\usepackage{enumitem}
\usepackage{multicol}
\usepackage{venndiagram}

\usepackage{fancyhdr}
\pagestyle{fancy}
\fancyhf{}
\renewcommand{\headrulewidth}{0pt} % disable the underline of the header
\raggedbottom
\hfuzz=2mm %suppresses overfull box warnings

\usepackage{hyperref}

\fancyhead[LO]{La Scuola d'Italia / Dr.\ Huson / IB Math 12 \\* 3 November 2026}
\fancyhead[RO]{\fbox{\begin{minipage}{6cm}\footnotesize Nickname:\\[0.5cm]\end{minipage}}}
\fancyfoot[C]{\thepage}

\begin{document}
\subsubsection*{3.2 Classwork: Transformations of Trigonometric Functions}

All angles are in radians. You may use your GDC.

\begin{enumerate}[itemsep=0.15cm]

%--- Problem 1: f(x) = 4sin(2(x - pi/6)) - 1: describe the transformations from y = sin x; amplitude, period, range; sketch with exact max and min [10 marks] ---
%--- Source: Original (AA SL 2.11, 3.7) ---
\item[1.]

The function $f$ is defined by
$f(x) = 4\sin\!\left(2\left(x - \dfrac{\pi}{6}\right)\right) - 1$.

\begin{enumerate}[itemsep=0.1cm]
  \item Describe fully the sequence of transformations that maps the graph
  of $y = \sin x$ onto the graph of $y = f(x)$.
  \item Write down the amplitude, the period and the range of $f$.
  \item On the grid, sketch the graph of $y = f(x)$ for $0 \le x \le \pi$.
  Mark the maximum point and the minimum point, and write their
  \emph{exact} coordinates.
\end{enumerate}

\begin{tikzpicture}
  \draw (-0.6,0) rectangle (14.6,16.0);
  \foreach \y in {10,11,...,15} \draw[dotted] (0,\y)--(14.1,\y);
  % sketch grid: 1 cm per pi/12 across, 0.8 cm per unit up; origin at (0.8, 5.1)
  \begin{scope}[shift={(0.8,5.1)}]
    \draw[thin, color=gray!60, xstep=1cm, ystep=0.8cm] (0,-4) grid (12,3.2);
    \draw[thick, ->] (-0.3,0) -- (12.6,0) node[right] {$x$};
    \draw[thick, ->] (0,-4.3) -- (0,3.7) node[right] {$y$};
    \foreach \v in {-5,-4,...,4}
      \draw (2pt,0.8*\v) -- (-2pt,0.8*\v) node[left, font=\footnotesize] {$\v$};
    \foreach \k/\lab in {2/\frac{\pi}{6}, 4/\frac{\pi}{3}, 6/\frac{\pi}{2},
        8/\frac{2\pi}{3}, 10/\frac{5\pi}{6}, 12/\pi}
      \node[below, font=\footnotesize] at (\k,-4.0) {$\displaystyle\lab$};
    % solution curve (MS check): max (5pi/12, 3), min (11pi/12, -5)
    %\draw[thick, domain=0:12, samples=120, smooth]
    %  plot (\x, {0.8*(4*sin(deg(2*(\x*pi/12 - pi/6)))) - 0.8});
  \end{scope}
\end{tikzpicture}

\newpage

%--- Problem 2: Daylight hours in New York, D(t) = a cos(b(t - c)) + d; find a, d, b, c from max and min; evaluate today; GDC solve D > 14 [11 marks] ---
%--- Source: Original (AA SL 3.7, 2.11); data rounded from NYC sunrise/sunset tables ---
\item[2.]

The number of hours of daylight in New York, $D$, can be modelled by the
function
\[
  D(t) = a\cos\bigl(b(t - c)\bigr) + d, \qquad 0 \le t \le 12,
\]
where $t$ is the time in months after 1 January, $a, b, c, d > 0$ and
$c < 12$.

The longest day has $15.1$ hours of daylight, at $t = 5.7$ (21 June).
The shortest day has $9.3$ hours of daylight, at $t = 11.7$
(21 December).

\begin{enumerate}[itemsep=0.1cm]
  \item Find the value of $a$ and the value of $d$.
  \item Find the value of $b$ and the value of $c$.
  \item Today is 3 November, $t = 10.1$. Use the model to find the number
  of hours of daylight today.
  \item Find the length of time during the year, in months, for which
  New York has more than $14$ hours of daylight.
\end{enumerate}

\begin{tikzpicture}
  \draw (-0.6,0) rectangle (14.6,14.0);
  \foreach \y in {1,2,...,13} \draw[dotted] (0,\y)--(14.1,\y);
\end{tikzpicture}

\end{enumerate}
\end{document}
